\documentclass{article}
\usepackage[utf8]{inputenc}
\usepackage[T1]{fontenc}
\usepackage{geometry}
\usepackage{graphicx}
\usepackage{float}
\usepackage{hyperref}
\usepackage[english,indonesian]{babel}

\title{Optimasi Pemrosesan Template LaTeX Menggunakan Pendekatan Berbasis AI}
\author{Academic Support Agent}
\date{\today}

\begin{document}
\maketitle

\begin{abstract}

\end{abstract}

\chapter{PENDAHULUAN}


\chapter{TINJAUAN PUSTAKA}


% Original content (commented out):
% \documentclass[11pt, Spanish]{article}
\usepackage{graphicx}
\usepackage[colorlinks=true, linkcolor=blue]{hyperref}
\usepackage[spanish]{babel}
\selectlanguage{spanish}
\usepackage[utf8]{inputenc}
\usepackage[svgnames]{xcolor}



\usepackage{listings}
\usepackage{afterpage}
\pagestyle{plain}

\definecolor{dkgreen}{rgb}{0,0.6,0}
\definecolor{gray}{rgb}{0.5,0.5,0.5}
\definecolor{mauve}{rgb}{0.58,0,0.82}

%\lstset{language=R,
%    basicstyle=\small\ttfamily,
%   stringstyle=\color{DarkGreen},
%    otherkeywords={0,1,2,3,4,5,6,7,8,9},
%    morekeywords={TRUE,FALSE},
%    deletekeywords={data,frame,length,as,character},
%    keywordstyle=\color{blue},
%    commentstyle=\color{DarkGreen},
%}

\lstset{frame=tb,
language=R,
aboveskip=3mm,
belowskip=3mm,
showstringspaces=false,
columns=flexible,
numbers=none,
keywordstyle=\color{blue},
numberstyle=\tiny\color{gray},
commentstyle=\color{dkgreen},
stringstyle=\color{mauve},
breaklines=true,
breakatwhitespace=true,
tabsize=3
}

\usepackage{here}


\textheight=21cm
\textwidth=17cm
%\topmargin=-1cm
\oddsidemargin=0cm
\parindent=0mm
\pagestyle{plain}

%%%%%%%%%%%%%%%%%%%%%%%%%%
% La siguiente instrucción pone el curso automáticamente%
%%%%%%%%%%%%%%%%%%%%%%%%%%

\usepackage{color}
\usepackage{ragged2e}

\global\let\date\relax
\newcounter{unomenos}
\setcounter{unomenos}{\number\year}
\addtocounter{unomenos}{-1}
\stepcounter{unomenos}
\gdef\@date{ Primer Semestre \arabic{unomenos} }

\begin{document}

\begin{titlepage}

\begin{center}
\vspace*{-1in}

\begin{figure}[htb]
\begin{center}
\includegraphics[width=8cm]{logo_utp}
\end{center}
\end{figure}

\vspace*{0.30in}

UNIVERSIDAD TECNOLÓGICA DE PANAMÁ - \@date\\
\vspace*{0.30in}
FACULTAD DE INGIENERÍA ELÉCTRICA \\
\vspace*{0.30in}
\begin{large}
DISEÑO DE LINEAS Y SUBESTACIONES\\
\end{large}
\vspace*{0.2in}
\begin{Large}
\textbf{Informe de gira a Subestación David} \\
\end{Large}
\vspace*{0.3in}
\begin{large}
Profesor Edwin Aparicio \\
\end{large}
\vspace*{0.1in}
\rule{80mm}{0.1mm}\\
\vspace*{0.5in}
\begin{large}
Hecho por: \\
Jaime Caballero \\
\end{large}
\vspace*{0.5in}
\includegraphics[width=10cm]{fie_logo.png}
\end{center}
\end{titlepage}

\newcommand{\CC}{C\nolinebreak\hspace{-.05em}\raisebox{.4ex}{\tiny\bf +}\nolinebreak\hspace{-.10em}\raisebox{.4ex}{\tiny\bf +}}
\def\CC{{C\nolinebreak[4]\hspace{-.05em}\raisebox{.4ex}{\tiny\bf ++}}}

\tableofcontents
\newpage
\section{Introducción}

Las sociedades modernas requieren de grandes cantidades de grandes cantidades de energía eléctrica para poder funcionar. Esta energía generalmente es producida en sectores lejanos y puntuales de donde va a ser mayormente consumida.  La energía eléctrica para atravesar grandes distancias con la menor cantidad de perdidas tiene que ser elevada a muy altas tensiones y luego bajar de nuevo la tensión para poder ser consumida. El hecho que se tengan que dar diferentes niveles de voltaje, la caída de este debido a perdidas y lo compleja que es la carga (variante); hace que se necesiten lugares en donde efectuar estas operaciones y estas son las subestaciones eléctricas.  \\

La función principal de las subestaciones eléctricas es elevar o disminuir la tensión mediante el uso de transformadores. No obstante, también se utilizan para regular pequeñas fluctuaciones de voltaje debido a las perdidas, distribuir la energía a diferentes sectores, hacer mediciones importantes de diferentes parámetros de la red, aplicar medidas de protección contra fallas en la red. \\

Para realizar todas estas funciones las subestaciones requieren de equipos como interruptores, banco de baterías, seccionadoras, barras, transformadores de potencia y de medición. Todos estos equipos fueron examinados en la gira realizada a la subestación de David el día 10 de mayo de 2019 y se discutirán más a fondo a continuación. 

\section{Objetivos}

\begin{itemize}

    \item Reconocer la función e importancia de una subestación en un sistema de potencia. 
    \item Examinar los elementos principales de la subestación David y su función.  
    \item Enunciar las características de los equipos para comprender su dimensionamiento y propósito.  

\end{itemize}

\section{Descripción de la Subestación Eléctrica David }

La subestación de David es de transformación MT/MT (media tensión a media tensión), opera a una frecuencia de 60 Hz. Es de exterior AIS (Air Insulated Substation). Cuenta con una configuración de Barra sencilla-Interruptor amarre. Utiliza un sistema de protección de apoyo de recierre, este abre y cierra la línea al detectarse una falla hasta dos veces de permanecer totalmente abierto.  \\

%corregir ultima linea

\newpage

\begin{figure}[htb]
\begin{center}
\includegraphics[width=12cm]{unifilar_david.png}
\end{center}
\caption{Unifilar de la subestación David}
\end{figure}

A la subestación David entran dos líneas con un voltaje nominal de 34.5 kV, una proveniente de Mata del Nance y la otra de Valbuena. De esta salen 6 líneas con un voltaje nominal de 13.8 kV. \\

\begin{figure}[htb]
\begin{center}
\includegraphics[width=4cm]{Subestacion_david.png}
\end{center}
\caption{Subestación David}
\end{figure}



\newpage

\section{Equipos}

Una subestación necesita de equipos o elementos para poder realizar todas tareas de transformación, regulación, distribución y protección. Estos elementos deben ser dimensionados para los niveles de caracteristicas eléctricas que debe manejar la subestación. A continuación examinaremos los equipos de la subestación de David.


\subsection{Transformadores de Potencia}

La subestación David cuenta con 3 transformadores de Potencia, tienen una conexión Delta en el lado de alta y Estrella en el lado de Baja. Una relación de 34.5kV/13.8 KV. Cada uno cuenta con una capacidad de 14 MVA. Sus corrientes nominales son de 234 A en el lado de alta y 585 A. En el lado de baja tiene un sistema de enfriamiento ONAF (Inmerso en liquido con punto de inflamación menor o igual a 300 C, Convección natural a través del equipo de enfriamiento y ventilación, enfriamiento por aire, circulación forzada). Cuentan un TapChanger tipo HV-OFTC/LV-OLTC Además, cuenta con 3 transformadores para servicios auxiliares con una conexión Estrella-Estrella y 10MVA cada uno. 

% corregir la potencia de los transformadores auxiliares

\begin{table}[h]
  \begin{center}
    \label{tab:table1}
    \begin{tabular}{|c|c|} 
    \hline
      Potencia Nominal & 14 MVA \\
      \hline
     Voltaje máximo de Alta & 34.5 kV \\ \hline
     Voltaje máximo de Baja  & 13.8 kV \\ \hline
     Tipo de TAP & 16 posiciones \\ \hline
     Tipo de enfriamiento & ONAF \\ \hline
     Frecuencia & 60 Hz \\ \hline
  
    \end{tabular}
    \caption{Especificaciones técnicas del transformador de potencia}
  \end{center}
\end{table}

\begin{figure}[htb]
\begin{center}
\includegraphics[width=12cm]{transformador_potencia.png}
\end{center}
\caption{Transformador de potencia y su placa}
\end{figure}



\subsection{Interruptores}

Un interruptor es un dispositivo cuya función es interrumpir y restablecer la continuidad en un circuito eléctrico. La subestación David tiene 5 interruptores de protección para exteriores en el lado de alta tensión. Son de tipo: Tanque muerto, cuentan con un medio interruptivo de gas SF6, operan en un voltaje nominal de 34.5 kV, y una corriente nominal de 600 A. Su mecanismo de operación es resorte. 

\begin{table}[h]
  \begin{center}
    \label{tab:table1}
    \begin{tabular}{|c|c|} 
    \hline
      Medio Interruptivo & SF6 \\
      \hline
     Tipo de interruptor & Muerto \\ \hline
     Frecuencia del sistema  & 60 Hz \\ \hline
     Voltaje normal del sistema & 34.5 kV \\ \hline
     Corriente nominal & 600 A \\ \hline
     Corriente de corto circuito & 31500 A \\ \hline
     Mecanismo de operación & Tipo resorte \\ \hline
  
     
    \end{tabular}
    \caption{Especificaciones técnicas del interruptor}
  \end{center}
\end{table}

\begin{figure}[htb]
\begin{center}
\includegraphics[width=12cm]{interruptor_david.png}
\end{center}
\caption{Interruptores y  placa}
\end{figure}




\subsection{Seccionadoras}

Es un elemento que sirve para desconectar físicamente un circuito eléctrico. Por lo general se operan sin carga. La subestación tiene seccionadores monopolares. Cuenta con cuchillas manuelas y grupales. \\

\begin{table}[h]
  \begin{center}
    \label{tab:table1}
    \begin{tabular}{|c|c|} 
    \hline
      Voltaje Nominal & 34.5 kV \\
      \hline
     Tipo de accionamiento & Mono polar y tripolar \\ \hline
     Frecuencia  & 60 Hz \\ \hline
     Corriente nominal & 600 A \\ \hline
     Corriente Momentánea & 1200 A \\ \hline
    
    \end{tabular}
    \caption{Especificaciones técnicas de la Seccionadora }
  \end{center}
\end{table}


\subsection{Pararrayos}

Elemento que nos permite proteger la subestación contra sobretensiones causas por fenómenos atmosféricos. La subestación cuenta con 8 pararrayos. \\

\begin{table}[h]
  \begin{center}
    \label{tab:table1}
    \begin{tabular}{|c|c|} 
    \hline
      Tipo de Pararrayos & Silicio \\
      \hline
     Voltaje Nominal & 34.5 kV \\ \hline
     Clase de Pararrayo  & Punta Franklin \\ \hline
    \end{tabular}
    \caption{Especificaciones técnicas de los pararrayos}
  \end{center}
\end{table}


\subsection{Transformadores de corriente (\textit{CT})}

Los transformadores de corriente los utiliza la subestación para poder medir las altas corrientes que circulan por las líneas. Un transforman de corriente reduce la corriente que circula en las líneas mediante una relación de vueltas y así se puede medir una corriente relativamente baja. Las características del CT de la subestación son presentadas a continuación.

\begin{table}[h]
  \begin{center}
    \label{tab:table1}
    \begin{tabular}{|c|c|} 
    \hline
      Voltaje Nominal & 34.5 kV \\
      \hline
     Frecuencia de Operación & 60 Hz \\ \hline
     Corriente secundario  & 5 A \\ \hline
     Tipo de Multirrelación  & 2000:5 \\ \hline
    \end{tabular}
    \caption{Especificaciones técnicas del transformador de corriente}
  \end{center}
\end{table}

\begin{figure}[htb]
\begin{center}
\includegraphics[width=12cm]{transformador_ct.png}
\end{center}
\caption{Transformador de corriente y  placa}
\end{figure}

\subsection{Transformadores de potencial (\textit{PT})}

Los transformadores de potencial transforman el potencial, según una relación de vueltas, a un nivel bajo que es capaz de ser medido facilmente. Las características del PT de la subestación son:

\begin{table}[h]
  \begin{center}
    \label{tab:table1}
    \begin{tabular}{|c|c|} 
    \hline
      Tipo & monofásico \\
      \hline
     Conexión & Estrella \\ \hline
     Relación en alta  & 175/1 \\ \hline
      Relación en baja & 70/1 \\ \hline
    \end{tabular}
    \caption{Especificaciones técnicas del transformador de potencial}
  \end{center}
\end{table}

\subsection{Barras}

Estas son conjunto de conductores eléctricos que efectuan una conexión común entre los diferentes circuitos que consta una subestación. Las barras de la subestación  tienen una ampacidad de 2000 A. Cuentan con un tipo de aleación 6063-T6.
Tienen un diámetro de 4 pulgadas, y un diámetro por efecto corona de 1.5 pulgadas. \\

\begin{table}[h]
  \begin{center}
    \label{tab:table1}
    \begin{tabular}{|c|c|} 
    \hline
      Ampacidad & 2000 A \\
      \hline
     Tipo de Aleación & 6063-T6 \\ \hline
     Diámetro (in)  & 4 \\ \hline
     Diámetro por efecto corono (in)  & 1.5 \\ \hline
    \end{tabular}
    \caption{Especificaciones técnicas de la barra}
  \end{center}
\end{table}


\subsection{Banco de Capacitores}

Los bancos de capacitores contribuyen a corregir el factor de potencia y mejorar el perfil de voltaje, tanto en sistemas de baja como de media y alta tensión. La subestación David cuenta con un banco de capacitores con potencia de 1000 kVA
por fase. En total cuentan con 3000 kVA. \\

\begin{figure}[htb]
\begin{center}
\includegraphics[width=4cm]{banco_capacitores.png}
\end{center}
\caption{Banco de capacitores}
\end{figure}

\newpage

\subsection{Banco de baterías}

Estos son los responsables de suministrar voltajes de baja tensión para los sistemas de protección, control, iluminación; especialmente cuando no hay energía debido a fallo. El banco de la subestación cuenta con las siguientes características.

\begin{table}[h]
  \begin{center}
    \label{tab:table1}
    \begin{tabular}{|c|c|} 
    \hline
      Voltaje & 132 V(dc) \\
      \hline
     Cargador & 35 A \\ \hline
     Duración  & 4 h \\ \hline
 
  
     
    \end{tabular}
    \caption{Características del banco de baterías}
  \end{center}
\end{table}

\begin{figure}[htb]
\begin{center}
\includegraphics[width=4cm]{banco_baterias.png}
\end{center}
\caption{Banco de baterías}
\end{figure}

\newpage

\section{Conclusión}


Una subestación es una parte fundamental de un sistema de potencia. Debido a los diferentes niveles de voltaje, distribución, las pérdidas y las fallas que se pueden dar; es necesario contar pon subestaciones en ciertos puntos de una red de potencia. Para que estas cumplan con todas estas funciones, necesitan de varios elementos que trabajan en conjunto. En la subestación de David pudimos observarlos y examinarlos. No solo comprender su función de la mano de un ingeniero, pero también su dimensionamiento de acuerdo a los niveles de tensión que va a manejar.  \\

En conclusión, para un sistema de potencia es completamente necesario el uso de subestaciones. Existen de diferentes tipos según su función específica, nuestra visita a la subestación de transformación de MT/MT de David nos permitió aprender más fondo acerca de los elementos que una subestación y como función en conjunto para suplir de energía de un país. 

%uoooooooooooooooo tumadreuooooooooooooooooooo UOOOOOOOOOOOOOOOOOOOOOOOOOOOOOOOOOOOOOOOOO
%AL FIN SE TERMINA ESTA PUTA MIERDA!!!!
%USEGREAS OSTOJEOGIRN ojeogiek


\end{document}

\end{document}